% !TeX root = ../navier-stokes-manuscript.tex

A first singular time would have to end a velocity field that is smooth on
every earlier cutoff while leaving no terminal state from which the same
Navier--Stokes history can restart.  For the maximal classical solution on
\([0,T_*)\), we locate the first possible failure in the adjacent trace pair
\(u\in L^2(0,T_*;H^2)\),
\(u_t\in L^2(0,T_*;L^2)\).  Its cutoff-uniform bound is equivalent to
\[
  D_*:=\int_0^{T_*}\|\Lambda^{3/2}u(t)\|_2^2\,dt<\infty,
\]
and the critical-height balance makes the same condition equivalent to an
upper bound on the signed critical-production prefix.  Once \(D_*\) is finite,
the enstrophy equation supplies a uniform \(H^1\) trace; the
Escauriaza--Seregin--\v{S}ver\'ak endpoint criterion then excludes a finite
\(T_*\).  If the corresponding adjacent estimate is available at every finite
temporal and spatial order, the rectangles contain identical derivatives on
their overlaps.  Their compatible intersection reaches one
\(u_*\in\bigcap_{m<\infty}H^m(\R^3)\), recovers the normalized pressure-complete
jet there, and restarts the same history.  This proves global smoothness
conditional on the whole-terminal rectangle condition.

The argument also identifies the boundary of that condition.  Viscosity pays
one higher spatial row on every fixed interval \(T<T_*\); small critical data
absorb the nonlinear production; and any hypothetical finite maximal time
obeys the datum-dependent capture horizon
\(T_*\leq C^4\|u_0\|_2^4/(2\nu^5)\).  Kinetic energy pays the contribution of
every fixed frequency range, so any failure of \(D_*\) is ultraviolet.  For an
arbitrary large Schwartz datum, however, the classical energy and enstrophy
balances do not provide the required upper bound on \(D_*\).  The manuscript
thus gives an exact endpoint reduction and its conditional same-history
closure, not an unconditional solution of the three-dimensional regularity
problem.
